% Vietnamese translation for OI Public Library
% Source-File: IOI中国国家候选队论文/2007/day1/7.袁昕颢《动态树问题及其应用》.ppt
\documentclass[11pt]{article}
\usepackage{oipl-vn}

\newcommand{\Oh}{\mathcal{O}}

\title{Bản trình chiếu: Bài toán cây động và ứng dụng}
\author{Yuan Xinhao}
\date{2007}

\begin{document}
\maketitle

\origtitle{Slide companion for dynamic trees and applications}
\sourcepath{IOI China candidate team papers / 2007 / day 1 / Yuan Xinhao.ppt}

\section{Phạm vi bản dịch}

Tệp gốc chỉ có PowerPoint và phần chữ trích xuất được rất hạn chế. Tên bài
và các dấu hiệu trong tệp cho thấy nội dung là cây động, tức duy trì một rừng
thay đổi theo thời gian dưới các thao tác nối, cắt và truy vấn đường đi.

\section{Bài toán cây động}

Trong nhiều bài đồ thị, cây không cố định: cạnh có thể được thêm vào hoặc xóa
đi, trong khi ta vẫn cần trả lời truy vấn về tổ tiên, đại diện thành phần,
giá trị cực trị trên đường đi, hoặc quan hệ liên thông. Nếu sau mỗi thay đổi
ta dựng lại toàn bộ cây, chi phí thường quá cao.

\section{Các thao tác cốt lõi}

Một cấu trúc cây động thường cần hỗ trợ:
\begin{itemize}
  \item \textsc{Link}: nối hai cây bằng một cạnh mới;
  \item \textsc{Cut}: xóa một cạnh đang có trong rừng;
  \item \textsc{FindRoot}: tìm gốc hoặc đại diện của cây chứa một đỉnh;
  \item \textsc{PathQuery}: truy vấn thông tin trên đường đi giữa hai đỉnh.
\end{itemize}

Link-cut tree là công cụ kinh điển cho nhóm thao tác này. Nó biểu diễn các
đường ưu tiên bằng splay tree, nhờ đó mỗi thao tác có chi phí khấu hao
\(\Oh(\log n)\).

\section{Ứng dụng}

Cây động thường xuất hiện trong bài toán cây khung động, luồng hoặc cắt trên
cây phụ trợ, duy trì đường đi nặng, và các bài cần thay đổi quan hệ cha con
trực tuyến. Khi cạnh mang trọng số, mỗi nút trong cấu trúc phụ có thể lưu cực
đại, cực tiểu hoặc tổng trên đường để trả lời truy vấn sau khi đưa đường cần
xét về một splay.

\section{Ghi chú}

Do slide gốc chủ yếu là hình và hiệu ứng, bản ghi chú này chỉ ghi lại khung
kiến thức chắc chắn thuộc chủ đề. Khi tiếp tục hoàn thiện, nên đối chiếu với
bản trình chiếu gốc bằng công cụ đọc PowerPoint chuyên dụng để khôi phục các
ví dụ cụ thể.

\end{document}
